\section{Conclusion}

\label{sec:conclusion}
%% ============================================================================

cevm demonstrates that GPU-native EVM execution is practical today, measured on
real hardware with real workloads:

\begin{enumerate}
\item \textbf{GPU EVM kernel}: 85.9\,M opcodes/sec, 4.85$\times$ over CPU. GPU wins
      above 500 opcodes per transaction.
\item \textbf{C++ interpreter}: 20.9\,Ggas/s with full state, 5.1$\times$ over Go EVM.
\item \textbf{Parallel ecrecover}: 7.1$\times$ on 10 cores. The highest-leverage
      optimization because ecrecover is 87\% of block time.
\item \textbf{Post-quantum crypto}: All three NIST FIPS algorithms on GPU via shared
      NTT primitive.
\item \textbf{FHE}: TFHE blind rotation and external product on GPU for encrypted
      smart contract execution.
\end{enumerate}

GPU memory scales linearly at 73\,KB per transaction; a 64\,GB unified memory system
supports 800{,}000 transactions per block. The system is validated by 1{,}183 tests
across five languages and three shader platforms.

All 90 shader files (30 Metal, 30 CUDA, 30 WGSL) are complete with zero stubs remaining.
Montgomery constants are verified byte-identical across Metal and CUDA backends; Keccak-256
round constants are verified against FIPS~202. 12 of 18 EVM precompiles are GPU-accelerated.
GPU signature recovery achieves 32$\times$ speedup (1613\,ms $\to$ $\sim$50\,ms for 47K
signatures), and the consensus pipeline delivers 107K blocks/sec on M1 Max with a 100M TPS
theoretical ceiling via BurstParams. Red team audit: 0 critical, 0 high findings---approved
for ship. CUDA/H100 ports are ready for CI validation and projected to deliver 4--6$\times$
additional throughput over the current Metal implementation.

\begin{thebibliography}{9}

\bibitem{blockstm}
A.~Gelashvili, A.~Spiegelman, Z.~Xiang, G.~Danezis, Z.~Li, Y.~Malkhi, Y.~Xia, R.~Zhou.
\emph{Block-STM: Scaling Blockchain Execution by Turning Ordering Curse to a Performance Blessing}.
PPoPP 2023.

\bibitem{pevm}
Risechain.
\emph{pevm: Parallel EVM implementation in Rust}.
\url{https://github.com/risechain/pevm}, 2024.

\bibitem{reth}
Paradigm.
\emph{reth: Modular, contributor-friendly and blazing-fast implementation of the Ethereum protocol, in Rust}.
\url{https://github.com/paradigmxyz/reth}, 2023.

\bibitem{evmone}
ipsilon.
\emph{evmone: Fast Ethereum Virtual Machine implementation}.
\url{https://github.com/ipsilon/evmone}, 2019.

\bibitem{zap}
Hanzo AI.
\emph{ZAP: Zero-Copy Application Protocol}.
\url{https://github.com/zap-proto/zap}, 2025.

\bibitem{quasar}
Lux Industries, Inc.
\emph{Quasar: Post-Quantum Finality Consensus for Multi-Chain Networks}.
Lux Technical Report, 2025.

\bibitem{tfhe}
I.~Chillotti, N.~Gama, M.~Georgieva, M.~Izabach\`{e}ne.
\emph{TFHE: Fast Fully Homomorphic Encryption over the Torus}.
Journal of Cryptology, 33(1), 2020.

\bibitem{mldsa}
NIST.
\emph{FIPS 204: Module-Lattice-Based Digital Signature Standard (ML-DSA)}.
National Institute of Standards and Technology, 2024.

\bibitem{mlkem}
NIST.
\emph{FIPS 203: Module-Lattice-Based Key-Encapsulation Mechanism Standard (ML-KEM)}.
National Institute of Standards and Technology, 2024.

\bibitem{slhdsa}
NIST.
\emph{FIPS 205: Stateless Hash-Based Digital Signature Standard (SLH-DSA)}.
National Institute of Standards and Technology, 2024.

\bibitem{cuevm}
FaceBlockTeam.
\emph{CuEVM: CUDA-based Ethereum Virtual Machine}.
\url{https://github.com/FaceBlockTeam/gpuEVM}, 2024.

\bibitem{gatlingx}
GatlingX.
\emph{GatlingX: GPU-Accelerated EVM Execution}.
2024. No reproducible benchmarks available.

\bibitem{parallelevm}
Y.~Chen, S.~Lu, J.~Fan, Z.~Li, Z.~Liu.
\emph{ParallelEVM: Leveraging Hardware Parallelism for Ethereum Virtual Machine Execution}.
EuroSys 2025.

\bibitem{gpumpt}
L.~Wu, Z.~Liu, Y.~Chen.
\emph{GPU-MPT: GPU-Accelerated Merkle Patricia Trie for Blockchain State}.
VLDB 2024.

\end{thebibliography}

